\documentclass[12pt]{article}
\usepackage[utf8]{inputenc}
\usepackage[T1]{fontenc}
\usepackage{hyperref}
\title{Southern Tujia Language Resources for K-12 Curriculum Development}
\author{Compiled by Claude Sonnet 5 (claude-sonnet-5) for Andr\'as Kornai}
\date{}
\begin{document}
\maketitle
\section{Introduction}
Southern Tujia (ISO 639-3 tjs) is a Sino-Tibetan language of Hunan, China. It is not mutually intelligible with Northern Tujia (tji). Search results place the remaining speakers in a few hamlets of Tanxi Township, Luxi County, and give figures of about 1,500 speakers in 3 villages. Little is publicly available, and nothing found is a school curriculum.
\section{Primary Educational Resources}
No primers, textbooks or school programmes for Southern Tujia turned up in the searches.
\section{Linguistic and Reference Materials}
\begin{itemize}
\item Wikipedia, ``Tujia language'': general background. \url{https://en.wikipedia.org/wiki/Tujia_language}
\item MPI Language Survey, Southern Tujia numerals. \url{https://lingweb.eva.mpg.de/channumerals/Tujia-Southern.htm}
\item Endangered Languages Project entry for Southern Tujia. \url{https://www.endangeredlanguages.com/lang/1744?hl=de}
\end{itemize}
The search summaries mention Zhang (2006) on the Poluozhai dialect, Luxi County, and Brassett \& Brassett (2005). I did not see a URL for either, so they are not listed as resources.
\section{Religious and Community Resources}
None found.
\section{Digital and Audio Resources}
\begin{itemize}
\item ELDP/HRELP archive, ``Documentation of the Southern Tujia Language of China''. The search result showed only the record title, so I did not see its contents. \url{https://eldp.access.preservica.com/dk0260/}
\end{itemize}
\section{Curriculum Development Context}
Southern Tujia has very few speakers and, as far as the searches show, no teaching materials. A curriculum would have to be built from the archived documentation materials and from work with the remaining speakers. Whether the archive is open to use should be checked directly.
\end{document}